
\documentclass[12pt]{article}
\usepackage{amsmath}
\usepackage{amsfonts}
\usepackage{amssymb}
\usepackage{geometry}
\geometry{a4paper, margin=0.8in}
\pagenumbering{gobble}

\title{Homework solutions for 02YELMA \\ Chapter: Electrostatics}
\begin{document}
\date{}
\maketitle
\vspace{-0.8in}

\begin{enumerate}

\item At the center, each charge produces an electric field of equal magnitude. Due to symmetry, the contributions cancel pairwise:
\[
\vec{E}_{\text{net}} = 0.
\]

If one charge is replaced by $-q$, symmetry is broken. The three $+q$ charges and one $-q$ charge no longer cancel. The resulting field points toward the negative charge. By summing vector contributions, the magnitude is:
\[
E = \frac{1}{4\pi\varepsilon_0}\frac{2\sqrt{2}\,q}{a^2}.
\]

The force on a test charge $Q$ at the center is:
\[
\vec{F} = Q \vec{E}.
\]

\item By symmetry, only the axial ($z$) component survives. Each charge element contributes a field whose radial components cancel.

\[
E_z = \frac{1}{4\pi\varepsilon_0} \frac{z (2\pi r \lambda)}{(r^2 + z^2)^{3/2}}.
\]

Thus:
\[
\vec{E} = \frac{1}{4\pi\varepsilon_0}\frac{2\pi r \lambda z}{(r^2+z^2)^{3/2}} \hat{z}.
\]

\item Using Gauss’s law, we first compute the enclosed charge:
\[
Q_{\text{enc}}(r) = \int_0^r kr' \cdot 4\pi r'^2 dr' = \pi k r^4.
\]

Then:
\[
E(4\pi r^2) = \frac{Q_{\text{enc}}}{\varepsilon_0}
\Rightarrow
E_{\text{in}} = \frac{k r^2}{4\varepsilon_0}.
\]

Outside the sphere, the field behaves like a point charge:
\[
Q_{\text{tot}} = \pi k R^4, \quad
E_{\text{out}} = \frac{1}{4\pi\varepsilon_0}\frac{\pi k R^4}{r^2}.
\]

\item The potential is obtained by integrating the electric field:
\[
V(r) = -\int E \, dr.
\]

Outside the sphere:
\[
V(r) = \frac{1}{4\pi\varepsilon_0}\frac{\pi k R^4}{r}.
\]

Inside, the potential is continuous and found by integrating inward:
\[
V(r) = V(R) - \int_r^R E_{\text{in}} \, dr.
\]

\item Field lines radiate outward symmetrically from each corner and curve away from the center. The field is strongest near the corners and weakest at the center due to cancellation.

\item By integrating concentric rings:

\[
V(z) = \frac{\sigma}{2\varepsilon_0}\left(\sqrt{R^2 + z^2} - z\right).
\]

The electric field is obtained via:
\[
E_z = -\frac{dV}{dz} = \frac{\sigma}{2\varepsilon_0}\left(1 - \frac{z}{\sqrt{R^2+z^2}}\right).
\]

\item Total charge:
\[
Q = 4\pi \int_0^R \rho_0\left(1 - \frac{r}{R}\right) r^2 dr = \frac{\pi \rho_0 R^3}{3}.
\]

Using Gauss’s law:
\[
E_{\text{in}} \propto \frac{Q_{\text{enc}}}{r^2}, \quad
E_{\text{out}} = \frac{1}{4\pi\varepsilon_0}\frac{Q}{r^2}.
\]

Energy:
\[
U = \frac{\varepsilon_0}{2}\int E^2 d\tau.
\]

Final result:
\[
U = \frac{5}{7} \cdot \frac{3Q^2}{20\pi\varepsilon_0 R}.
\]

This is smaller than the uniform sphere energy, since charge is more concentrated near the center.

\item Energy is the sum over pair interactions:
\[
U = \frac{1}{4\pi\varepsilon_0}\sum_{i<j} \frac{q_i q_j}{r_{ij}}.
\]

For the triangle:
\[
U = \frac{1}{4\pi\varepsilon_0}\left(\frac{q^2}{a} - \frac{2q^2}{a} - \frac{2q^2}{a}\right)
= -\frac{3q^2}{4\pi\varepsilon_0 a}.
\]

For the centered configuration, the distances change, leading to a more negative energy, hence more stable.

\item 

\[
U = \frac{\varepsilon_0}{2} \int E^2 d\tau.
\]

For a point charge:
\[
E = \frac{1}{4\pi\varepsilon_0}\frac{q}{r^2}.
\]

The integral diverges at $r=0$, so:
\[
U = \infty.
\]

This shows that a point charge has infinite self-energy.

\item We model the electron as a spherical shell of radius $R$ with total charge $e$ uniformly distributed on its surface. Electric field of a charged shell by Gauss’s law:

- For $r > R$ (outside):
\[
E(r) = \frac{1}{4\pi \varepsilon_0} \frac{e}{r^2}.
\]

- For $r < R$ (inside):
\[
E(r) = 0.
\]

Thus, the electric field exists only outside the shell. The energy stored in the electric field is:
\[
U = \frac{\varepsilon_0}{2} \int_{\text{all space}} E^2 \, d\tau.
\]

Since $E=0$ inside the shell, we only integrate over $r \ge R$. Substitute the field:

\[
E^2 = \left(\frac{1}{4\pi \varepsilon_0} \frac{e}{r^2}\right)^2
= \frac{e^2}{16 \pi^2 \varepsilon_0^2 r^4}.
\]

The volume element in spherical coordinates:
\[
d\tau = 4\pi r^2 dr.
\]

Therefore, the energy integral becomes:

\[
U = \frac{\varepsilon_0}{2} \int_R^\infty E^2 \, d\tau
= \frac{\varepsilon_0}{2} \int_R^\infty \frac{e^2}{16 \pi^2 \varepsilon_0^2 r^4} \cdot 4\pi r^2 \, dr.
\]

Simplify:

\[
U = \frac{\varepsilon_0}{2} \cdot \frac{e^2}{16 \pi^2 \varepsilon_0^2} \cdot 4\pi \int_R^\infty \frac{1}{r^2} \, dr.
\]

\[
U = \frac{e^2}{8\pi \varepsilon_0} \int_R^\infty \frac{1}{r^2} \, dr.
\]

\[
\int_R^\infty \frac{1}{r^2} dr = \left[-\frac{1}{r}\right]_R^\infty = \frac{1}{R}.
\]

\[
U = \frac{e^2}{8\pi \varepsilon_0 R}.
\]

Equating this energy to the electron rest energy:
\[
mc^2 = \frac{e^2}{8\pi \varepsilon_0 R},
\]
Thus,
\[
R = \frac{1}{8\pi \varepsilon_0} \frac{e^2}{mc^2}.
\]

This is the standard classical estimate of the electron size, which is about $2.8 \times 10^{-15}$ meters.

\item Using the pressure expression on the surface of a conductor:
\[
F = \frac{\varepsilon_0}{2} E^2 A.
\]

For a charged sphere:
\[
E = \frac{Q}{4\pi\varepsilon_0 R^2}
\Rightarrow
F = \frac{Q^2}{32\pi^2 \varepsilon_0 R^2}.
\]

\item The pressure between plates:
\[
P = \frac{\varepsilon_0}{2}E^2.
\]

Work done:
\[
W = P A \epsilon = \frac{\varepsilon_0}{2}E^2 A \epsilon.
\]

\item Consider two coaxial cylinders of radii $a$ (inner) and $b$ (outer), with charge per unit length $\lambda$ on the inner cylinder and $-\lambda$ on the outer. Choose a cylindrical Gaussian surface of radius $r$ and length $L$ (with $a < r < b$):

\[
\oint \vec{E}\cdot d\vec{A} = E(2\pi r L) = \frac{\lambda L}{\varepsilon_0}.
\]

Thus,
\[
E(r) = \frac{\lambda}{2\pi \varepsilon_0 r}.
\]

The potential difference between the cylinders is:

\[
V = V(a) - V(b) = -\int_a^b \vec{E}\cdot d\vec{r}.
\]

Since $\vec{E}$ is radial:

\[
V = -\int_a^b \frac{\lambda}{2\pi \varepsilon_0 r} \, dr = -\frac{\lambda}{2\pi \varepsilon_0} \int_a^b \frac{1}{r} dr
= -\frac{\lambda}{2\pi \varepsilon_0} \left[\ln r\right]_a^b = -\frac{\lambda}{2\pi \varepsilon_0} \ln\left(\frac{b}{a}\right).
\]

Taking magnitude:

\[
V = \frac{\lambda}{2\pi \varepsilon_0} \ln\left(\frac{b}{a}\right).
\]

By definition:

\[
C' = \frac{\lambda}{V} = \frac{\lambda}{\frac{\lambda}{2\pi \varepsilon_0} \ln(b/a)} = \frac{2\pi \varepsilon_0}{\ln(b/a)}.
\]

\item We compute the multipole expansion of the potential at large distance $r \gg a$. Total charge:

\[
Q = \sum_i q_i = (-2q) + (-2q) + q + 3q = 0.
\]

Thus the monopole term vanishes:
\[
V_{\text{mono}} = \frac{1}{4\pi\varepsilon_0}\frac{Q}{r} = 0.
\]
 
Dipole moment by definition:

\[
\vec{p} = \sum_i q_i \vec{r}_i.
\]

Compute contributions:

\[
(-2q)(0,-a,0) = (0,2qa,0),
\]
\[
(-2q)(0,a,0) = (0,-2qa,0),
\]
\[
q(0,0,-a) = (0,0,-qa),
\]
\[
3q(0,0,a) = (0,0,3qa).
\]

Sum:
\[
\vec{p} = (0,0,2qa).
\]

So the dipole moment is:
\[
\vec{p} = 2qa\,\hat{z}.
\]

For $r \gg a$, the leading contribution is the dipole term:

\[
V(\vec{r}) \approx \frac{1}{4\pi\varepsilon_0} \frac{\vec{p}\cdot \hat{r}}{r^2}.
\]

Since $\vec{p} \parallel \hat{z}$:

\[
\vec{p}\cdot \hat{r} = p \cos\theta.
\]

\[
V(r,\theta) \approx \frac{1}{4\pi\varepsilon_0} \frac{2qa \cos\theta}{r^2}.
\]


\item 

(a) Charges: $3q$ at $(0,0,a)$ and $-q$ at $(0,0,0)$. Monopole moment:

\[
Q = 3q + (-q) = 2q.
\]
Dipole moment:

\[
\vec{p} = \sum_i q_i \vec{r}_i = 3q(0,0,a) + (-q)(0,0,0) = (0,0,3qa).
\]

\[
\vec{p} = 3qa\,\hat{z}.
\]

Potential (far field):

\[
V(r,\theta) \approx \frac{1}{4\pi\varepsilon_0}
\left(
\frac{Q}{r} + \frac{\vec{p}\cdot \hat{r}}{r^2}
\right).
\]

Since $\vec{p}\parallel \hat{z}$:
\[
\vec{p}\cdot \hat{r} = p\cos\theta.
\]

\[
V(r,\theta) \approx \frac{1}{4\pi\varepsilon_0}
\left(
\frac{2q}{r} + \frac{3qa\cos\theta}{r^2}
\right).
\]

(b) Charges: $3q$ at $(0,0,0)$ and $-q$ at $(0,0,-a)$. Monopole moment:

\[
Q = 3q + (-q) = 2q.
\]

Dipole moment:

\[
\vec{p} = 3q(0,0,0) + (-q)(0,0,-a).
\]

\[
(-q)(0,0,-a) = (0,0,qa),
\]

so:
\[
\vec{p} = (0,0,qa) = qa\,\hat{z}.
\]

Potential (far field):

\[
V(r,\theta) \approx \frac{1}{4\pi\varepsilon_0}
\left(
\frac{2q}{r} + \frac{qa\cos\theta}{r^2}
\right).
\]



\item A pure dipole $\vec{p}$ is located at the origin and oriented in the $-z$ direction:
\[
\vec{p} = -p\,\hat{z}.
\]

A point charge $q$ is located at $(a,0,0)$. The position $(a,0,0)$ lies on the $x$-axis. In spherical coordinates:
\[
r = a, \quad \theta = \frac{\pi}{2}.
\]

Given:
\[
\vec{E}(r,\theta) = \frac{p}{4\pi\varepsilon_0 r^3}
\left( 2\cos\theta\,\hat{r} + \sin\theta\,\hat{\theta} \right).
\]

Substitute $\theta = \frac{\pi}{2}$:
\[
\cos\theta = 0, \quad \sin\theta = 1.
\]

\[
\Rightarrow \vec{E} = \frac{p}{4\pi\varepsilon_0 a^3}\,\hat{\theta}.
\]

At $\theta = \frac{\pi}{2}$, the unit vector $\hat{\theta}$ points in the $-z$ direction:
\[
\hat{\theta} = -\hat{z}.
\]

Since the dipole is already oriented along $-z$, the field becomes:
\[
\vec{E} = -\frac{p}{4\pi\varepsilon_0 a^3}\,\hat{z}.
\]

Force on the charge:

\[
\vec{F} = q \vec{E}.
\]

\[
\Rightarrow \vec{F} = -\frac{qp}{4\pi\varepsilon_0 a^3}\,\hat{z}.
\]


\item Two dipoles are separated along the $x$-axis by a distance $r$:
\[
\vec{p}_1 = p_1 \hat{z}, \quad \vec{p}_2 = p_2 \hat{x}.
\]

We compute the torque on $\vec{p}_2$ due to the field created by $\vec{p}_1$. The field point (location of $\vec{p}_2$) lies on the $x$-axis. Therefore:
\[
\theta = \frac{\pi}{2}
\]
with respect to $\vec{p}_1$ (which points along $z$). Given:
\[
\vec{E}_1(r,\theta) = \frac{p_1}{4\pi\varepsilon_0 r^3}
\left( 2\cos\theta\,\hat{r} + \sin\theta\,\hat{\theta} \right).
\]

At $\theta = \frac{\pi}{2}$:
\[
\cos\theta = 0, \quad \sin\theta = 1,
\]

so:
\[
\vec{E}_1 = \frac{p_1}{4\pi\varepsilon_0 r^3}\,\hat{\theta}.
\]

At the point on the $x$-axis (equatorial plane of the dipole), $\hat{\theta}$ points in the $-z$ direction:
\[
\hat{\theta} = -\hat{z}.
\]

Thus:
\[
\vec{E}_1 = -\frac{p_1}{4\pi\varepsilon_0 r^3}\,\hat{z}.
\]

The torque on a dipole in an external field is:
\[
\vec{\tau} = \vec{p}_2 \times \vec{E}_1.
\]

Substitute:
\[
\vec{\tau} = (p_2 \hat{x}) \times \left(-\frac{p_1}{4\pi\varepsilon_0 r^3}\hat{z}\right).
\]

\[
\hat{x} \times \hat{z} = -\hat{y}.
\]

So:
\[
\vec{\tau} = -\frac{p_1 p_2}{4\pi\varepsilon_0 r^3} (\hat{x} \times \hat{z})
= -\frac{p_1 p_2}{4\pi\varepsilon_0 r^3} (-\hat{y}).
\]

\[
\vec{\tau} = \frac{p_1 p_2}{4\pi\varepsilon_0 r^3}\,\hat{y}.
\]


\end{enumerate}

\end{document}
